% Part: model-theory
% Chapter: interpolation
% Section: introduction

\documentclass[../../../include/open-logic-section]{subfiles}

\begin{document}

\olfileid{mod}{int}{int}

\olsection{تعارف}

درمیانی جملے دے قضیے دا بیان ایہ اے: فرض کرو دتے دو
فارمولے جملے نیں تے
$\Entails !A \lif !B$۔ تاں اک !!a{sentence} $!C$
موجود اے جس لئی $\Entails !A \lif !C$ تے
$\Entails !C \lif !B$۔ ایس توں ودھ، ہر
!!{constant}، !!{function} تے !!{predicate}
($\eq$ توں علاوہ) جو $!C$ وچ اے، $!A$
تے~$!B$ دوہاں وچ وی آؤندا اے۔
!!{sentence} $!C$ نوں
$!A$ تے~$!B$ دا \emph{درمیانی جملہ
(interpolant)} آکھیا جاندا اے۔

درمیانی جملے دا قضیہ اپنی تھاں دلچسپ اے، پر
اوہدی بنیادی اہمیت ایہ اے کہ ایس نال کسے نظریے
وچ قابلِ تعریف ہون بارے نتائج ثابت کیتے جا سکدے
نیں، تے ایہ وی پتہ لگایا جا سکدا اے کہ کنہاں
شرطاں تحت دو بے تضاد نظرییاں نوں ملاؤن توں
بے تضاد نظریہ بندا اے۔ پہلے نتیجے نوں بیتھ دا
قابلِ تعریف ہون دا قضیہ آکھدے نیں؛ دوجے نوں
رابنسن دا مشترکہ بے تضادی دا قضیہ۔

\end{document}
